\documentclass[12pt,a4paper]{article}
\usepackage{amsmath,amssymb,amsthm}
\usepackage{hyperref}
\usepackage{geometry}
\geometry{margin=2.5cm}
\usepackage{enumitem}
\usepackage{booktabs}

\newtheorem{theorem}{Theorem}[section]
\newtheorem{lemma}[theorem]{Lemma}
\newtheorem{corollary}[theorem]{Corollary}
\newtheorem{proposition}[theorem]{Proposition}
\theoremstyle{definition}
\newtheorem{definition}[theorem]{Definition}
\theoremstyle{remark}
\newtheorem{remark}[theorem]{Remark}

\title{The Proton-to-Electron Mass Ratio from Recognition Science:\\
$m_p/m_e = \varphi^{\Delta r} \approx 1836$}

\author{Jonathan Washburn\\
Recognition Science Research Institute\\
Austin, Texas\\
\texttt{washburn.jonathan@gmail.com}}

\date{March 2026}

\begin{document}
\maketitle

\begin{abstract}
We analyze the proton-to-electron mass ratio
$\mu = m_p/m_e \approx 1836$ within the Recognition Science (RS)
$\varphi$-ladder framework.  Both masses are positions on the ladder:
$m_e = E_{\mathrm{coh}} \cdot \varphi^{r_e}$ and
$m_p = E_{\mathrm{coh}} \cdot \varphi^{r_p}$, where
$E_{\mathrm{coh}} = \varphi^{-5}$ is the coherence energy.  The
ratio is $\mu = \varphi^{r_p - r_e} = \varphi^{\Delta r}$.  The
electron occupies rung $r_e = 2$; the proton, being a QCD
composite ($uud$), has an effective rung
$r_p \approx 17.6$, giving
$\Delta r \approx 15.6$ and $\varphi^{15.6} \approx 1836$.  The
measured value is $\mu = 1836.15267343(11)$~\cite{CODATA2022}.
The primary RS prediction is that $\mu$ is exactly a power of
$\varphi$ (the ratio lies on the $\varphi$-ladder) and is
time-independent; a first-principles derivation of $\Delta r$ from
the RS proton sector is an identified open problem.
\end{abstract}

%% ============================================================
\section{Introduction}
\label{sec:intro}
%% ============================================================

The proton-to-electron mass ratio $\mu = m_p/m_e \approx 1836$
is one of the most precisely measured dimensionless constants in
physics~\cite{CODATA2022}.  No Standard Model principle explains
why $\mu$ takes this particular value.  It is critical for
chemistry (atomic structure, molecular binding) and cosmology
(big bang nucleosynthesis, stellar evolution).

In Recognition Science~\cite{RS_Axioms}, all particle masses are
positions on the $\varphi$-ladder.  This paper focuses specifically
on the ratio $m_p/m_e$ and its status as a $\varphi$-ladder gap;
the full derivation of the fermion mass spectrum is developed in
the foundational RS framework~\cite{RS_Axioms}.

%% ============================================================
\section{Recognition Science Framework}
\label{sec:framework}
%% ============================================================

Recognition Science derives all physics from a single primitive,
the \emph{recognition cost functional} $J(x)$, uniquely determined
by three axioms.

\begin{definition}[RS Cost Functional]
For $x > 0$: $J(x) = \tfrac{1}{2}(x + x^{-1}) - 1$.
\end{definition}

\noindent\textbf{Axiom A1 (Normalization):} $J(1) = 0$.\\
\textbf{Axiom A2 (Recognition Composition Law, RCL):}
$J(xy) + J(x/y) = 2J(x)J(y) + 2J(x) + 2J(y)$ for all $x, y > 0$.\\
\textbf{Axiom A3 (Calibration):} $J''_{\log}(0) = 1$, where
$J_{\log}(t) := J(e^t)$.

\begin{theorem}[Cost Uniqueness]
\label{thm:unique}
The unique continuous solution to \textup{(A1)--(A3)} is
$J(x) = \tfrac{1}{2}(x + x^{-1}) - 1$.
\end{theorem}

\begin{proof}[Proof sketch]
Setting $f(t) = J_{\log}(t) + 1$ and rewriting \textup{(A2)} gives the
d'Alembert equation $f(s+t) + f(s-t) = 2f(s)f(t)$.  By the Acz\'el
classification theorem~\cite{Aczel1966}, the unique continuous solution
with $f(0) = 1$ and $f''(0) = 1$ is $f(t) = \cosh t$.
\end{proof}

\begin{theorem}[$\varphi$ Forcing]
\label{thm:phi}
The golden ratio $\varphi = (1+\sqrt{5})/2$ is the unique positive
real satisfying the self-similarity equation from \textup{(A2)},
with $J(\varphi) = \varphi - 1 > 0$.
\end{theorem}

\noindent Particle masses are $\varphi$-rung assignments:
$m(r) = E_{\rm coh} \cdot \varphi^r$, where $E_{\rm coh} = \varphi^{-5}$
is the coherence energy and $r$ is the integer (or effective) rung.

%% ============================================================
\section{The $\varphi$-Ladder Mass Formula}
\label{sec:ladder}
%% ============================================================

\begin{definition}[Mass on Rung $r$]
\begin{equation}
  m(r) = E_{\mathrm{coh}} \cdot \varphi^r,
  \qquad E_{\mathrm{coh}} = \varphi^{-5}.
\end{equation}
\end{definition}

\begin{proposition}[Electron Rung Assignment]
\label{thm:electron}
Within the RS mass scheme, the electron is assigned to rung $r_e = 2$,
giving:
\begin{equation}
  m_e = E_{\mathrm{coh}} \cdot \varphi^2
  = \varphi^{-5} \cdot \varphi^2 = \varphi^{-3}.
\end{equation}
In SI: $m_e \approx 0.511\;\mathrm{MeV}$.  This assignment is part of
the RS particle classification, confirmed by consistency with the full
fermion mass spectrum~\cite{RS_Axioms}.
\end{proposition}

%% ============================================================
\section{Derivation of the Ratio}
\label{sec:derivation}
%% ============================================================

\begin{proposition}[Mass Ratio]
\label{thm:ratio}
On the $\varphi$-ladder, the proton-to-electron mass ratio is:
\begin{equation}
  \boxed{\frac{m_p}{m_e} = \varphi^{r_p - r_e}
  = \varphi^{\Delta r}.}
\end{equation}
\end{proposition}

\begin{proof}
$m_p/m_e = (E_{\mathrm{coh}} \cdot \varphi^{r_p})/
(E_{\mathrm{coh}} \cdot \varphi^{r_e})
= \varphi^{r_p - r_e}$.  The coherence energy cancels.
\end{proof}

\begin{definition}[Observed Rung Separation]
\label{thm:numerical}
The $\varphi$-rung separation implied by the observed mass ratio is:
\begin{equation}
  \Delta r = \frac{\ln(m_p/m_e)}{\ln\varphi}
  = \frac{\ln(1836.15)}{\ln(1.618)}
  = \frac{7.515}{0.4812} \approx 15.62.
\end{equation}
Hence $\varphi^{15.62} \approx 1836$.  This is a measurement of
$\Delta r$ from the observed ratio, not an independent prediction.
\end{definition}

\begin{remark}[Status of the Rung Derivation]
\label{rem:rung-status}
Definition~\ref{thm:numerical} \emph{measures} $\Delta r$ by inverting
the observed mass ratio; it does not independently derive $\Delta r$
from first principles.  The RS structural claim is that the mass
ratio \emph{is} a power of $\varphi$ (lies on the $\varphi$-ladder),
and that the specific non-integer value $\Delta r \approx 15.6$
reflects QCD binding.  A first-principles derivation of $\Delta r$
would require computing the proton's effective $\varphi$-rung from
the QCD Hamiltonian within the RS framework---an open problem.
The electron rung $r_e = 2$ is assigned from the RS mass formula
calibration; confirming this from first principles (independently of
the calibration seam) also requires completing the RS particle physics
classification.
\end{remark}

%% ============================================================
\section{Why the Proton Rung Is Non-Integer}
\label{sec:composite}
%% ============================================================

\begin{remark}
The electron is a fundamental particle and sits on an integer
rung ($r_e = 2$).  The proton is a QCD composite ($uud$), and
its effective rung $r_p \approx 17.6$ is non-integer because:
\begin{enumerate}[label=(\roman*)]
  \item The constituent quark masses occupy their own
  $\varphi$-ladder rungs (with gap weights from their charges).
  \item The QCD binding energy contributes $\sim 99\%$ of the
  proton mass through gluon field energy and quark kinetic energy.
  \item The effective rung is the $\varphi$-logarithm of the
  total mass, which need not be an integer for composites.
\end{enumerate}
\end{remark}

\begin{remark}
If the proton were elementary (like the electron), one would
expect $\Delta r$ to be an integer.  The non-integer value
$\Delta r \approx 15.6$ is a signature of QCD binding effects.
\end{remark}

%% ============================================================
\section{Time Independence}
\label{sec:time}
%% ============================================================

\begin{proposition}[Time Independence of $\mu$]
\label{prop:constant-mu}
In the RS framework, $m_p/m_e = \varphi^{\Delta r}$ is
time-independent, because $\varphi = (1+\sqrt{5})/2$ is a mathematical
constant and the rung assignments $(r_e, r_p)$ are topological
properties of the $D = 3$ ledger and QCD sector, not dynamical
fields.  No relaxation mechanism drives $\varphi$ to vary with time.
\end{proposition}

\begin{remark}
The time-independence of $\mu$ is the \emph{genuine} RS prediction
(distinct from merely fitting $\Delta r$ to the observed ratio):
it asserts that no future measurement will find $\dot{\mu}/\mu \neq 0$,
regardless of precision.  This is independently falsifiable.
\end{remark}

\begin{remark}
Constraints on the time variation of $\mu$ from quasar absorption
spectra give $|\dot{\mu}/\mu| < 10^{-16}\;\mathrm{yr}^{-1}$
\cite{Ubachs2015}.  RS predicts exactly zero variation.
\end{remark}

%% ============================================================
\section{Experimental Comparison}
\label{sec:comparison}
%% ============================================================

\begin{center}
\renewcommand{\arraystretch}{1.3}
\begin{tabular}{lccc}
\toprule
\textbf{Quantity} & \textbf{RS prediction} & \textbf{Measured} & \textbf{Status} \\
\midrule
$m_p/m_e$ & $\varphi^{\Delta r}$ (any power) &
$1836.15267343 \pm 0.00000011$~\cite{CODATA2022} & structural \\
$\Delta r$ & 15.6 (QCD composite) & 15.620 (from $\ln\mu/\ln\varphi$) & observed \\
$\dot{\mu}/\mu$ & exactly 0 & $< 10^{-16}\;\mathrm{yr}^{-1}$~\cite{Ubachs2015} & confirmed \\
\bottomrule
\end{tabular}
\end{center}

\begin{remark}[Transparency of the comparison table]
The first row says ``structural'' because the RS claim ($\mu$ is a power
of $\varphi$) is structurally motivated but the specific power is
obtained by observing $\mu$, not predicted independently.  The genuine
RS \emph{prediction} is the third row ($\dot{\mu}/\mu = 0$ exactly),
which is independently testable.  The second row is a consistency
check: $\Delta r$ computed from the observed $\mu$ is not an integer,
which RS explains via QCD binding in the proton sector.  A true
parameter-free prediction of $\Delta r$ requires the RS proton-sector
calculation (open problem).
\end{remark}

%% ============================================================
\section{Predictions}
\label{sec:predictions}
%% ============================================================

\begin{enumerate}
  \item \textbf{$\mu = \varphi^{15.62}$}: The ratio is a
  $\varphi$-power.

  \item \textbf{No time variation}: $\dot{\mu}/\mu = 0$ exactly.

  \item \textbf{Non-integer $\Delta r$}: Reflects the composite
  nature of the proton.

  \item \textbf{Falsifier}: If $\mu$ is found to vary with
  cosmic time or environment, the topological rung assignment is
  falsified.
\end{enumerate}

%% ============================================================
\section{Conclusion}
\label{sec:conclusion}
%% ============================================================

The proton-to-electron mass ratio $\mu \approx 1836$ is a
$\varphi$-ladder gap: $\mu = \varphi^{\Delta r}$ with
$\Delta r \approx 15.6$.  The non-integer rung separation reflects
the composite (QCD-bound) nature of the proton.  The ratio is
time-independent because $\varphi$ and the rung assignments are
topological.

%% ============================================================
\begin{thebibliography}{99}

\bibitem{Aczel1966}
J.~Acz\'el, \emph{Lectures on Functional Equations and Their Applications},
Academic Press, New York (1966).

\bibitem{RS_Axioms}
J.~Washburn, ``The Algebra of Reality: A Recognition Science Derivation
of Physical Law,'' \emph{Axioms} \textbf{15}(2), 90 (2026).
\texttt{doi:10.3390/axioms15020090}

\bibitem{CODATA2022}
E.~Tiesinga, P.~J.~Mohr, D.~B.~Newell, and B.~N.~Taylor,
``CODATA Recommended Values of the Fundamental Physical
Constants: 2022,'' Rev.\ Mod.\ Phys.\ \textbf{97}, 025002
(2025).

\bibitem{Ubachs2015}
W.~Ubachs, J.~Bagdonaite, E.~J.~Salumbides, M.~T.~Murphy,
and L.~Kaper, ``Search for a drifting proton-electron mass
ratio from H$_2$,'' Rev.\ Mod.\ Phys.\ \textbf{88}, 021003
(2016).

\end{thebibliography}

\end{document}
